\chapter{Every fixed congruence depth has actual prime survivors}\label{ch:congruences}
\section{A theorem about actual primes}
Congruences can disprove a zero at a particular prime, but no fixed finite congruence depth excludes every prime. The following statement concerns the actual coefficients of $\Delta$, not a synthetic recurrence.

\begin{theorem}\label{thm:actualsurvivors}
For every integer $M\ge1$, a set of rational primes of positive natural density satisfies
\begin{equation}\label{eq:actualsurvivors}
p\equiv-1\pmod{5M},\qquad\tau(p)\equiv0\pmod M.
\end{equation}
At every such prime, simultaneously for every $r\ge1$,
\begin{equation}\label{eq:allrcongruence}
R(p^r)\equiv0\pmod M.
\end{equation}
\end{theorem}

Natural density here is relative to the rational primes. The proof exhibits a positive-density Chebotarev subset of the displayed set. The full set is itself a union of conjugacy classes in the same finite quotient, so also has positive density. No numerical density for the full set is calculated.

The theorem is a consequence of the classical Galois representations and Chebotarev, followed by a translation to the primitive local factor. It is not a new general density theorem. Its usefulness here is to delimit one proposed nonvanishing strategy exactly.

\section{Integral representations and complex conjugation}
For each prime $\ell$, the representation attached to $\Delta$ has the arithmetic-Frobenius normalization
\[
\rho_\ell:G_{\Q}\longrightarrow GL_2(\Q_\ell),
\]
is continuous and unramified outside $\ell$, and satisfies
\begin{equation}\label{eq:frobpoly}
\det(X-\rho_\ell(\Frob_p))=X^2-\tau(p)X+p^{11},
\qquad p\ne\ell.
\end{equation}
Its determinant is the eleventh power of the $\ell$-adic cyclotomic character. The original cohomological formulation uses a dual convention for Frobenius. The trace and determinant contract in this normalization is explicitly checked in the primary sources \cite{DeligneModular,SerreBourbaki}. Both conventions agree at complex conjugation for the argument below.

A continuous representation of the compact group $G_{\Q}$ has a stable $\Z_\ell$ lattice. To see this directly, take an initial lattice $L_0$. Compactness bounds all its translates inside $\ell^{-a}L_0$ for some $a$. The module
\[
L=\sum_{g\in G_{\Q}}\rho_\ell(g)L_0
\]
is contained in $\ell^{-a}L_0$ and contains $L_0$. It is finitely generated, since a submodule of a finite free $\Z_\ell$ module is finitely generated, and it is invariant by construction. Reductions modulo $\ell^e$ are therefore defined. No image-surjectivity theorem is required.

Let $c$ be complex conjugation. Then $c^2=1$ and the cyclotomic character sends $c$ to $-1$. Consequently $\det\rho_\ell(c)=-1$. In characteristic zero an involution is diagonalizable, with eigenvalues among $1,-1$. The determinant forces one of each, so
\begin{equation}\label{eq:complexconjtrace}
\Tr\rho_\ell(c)=0
\end{equation}
as an exact $\Z_\ell$ identity. This reasoning also covers $\ell=2$: the trace is computed before reduction. Diagonalizing an involution in characteristic $2$ would not justify the same step.

\section{The actual joint image}
Factor $M=\prod_{\ell^e\Vert M}\ell^e$ and form the single joint finite map
\begin{equation}\label{eq:jointimage}
\Psi:G_{\Q}\longrightarrow
\left(\prod_{\ell^e\Vert M}GL_2(\Z/\ell^e\Z)\right)
\times(\Z/5M\Z)^\times.
\end{equation}
Its components are the integral representation reductions and the cyclotomic character modulo $5M$. Let $G$ be its actual image. The kernel defines a finite Galois extension of $\Q$. The image need not be the whole displayed product. Relations between different representations and determinants are allowed.

Let $C$ be the conjugacy class of $\Psi(c)$ in $G$. Chebotarev gives unramified primes whose Frobenius lies in $C$, with density $|C|/|G|>0$ \cite{Chebotarev}. The trace in each representation component is zero by \eqref{eq:complexconjtrace}, and the cyclotomic component is $-1$. Thus every $\ell^e$ dividing $M$ divides $\tau(p)$, while $p\equiv-1\pmod{5M}$. This proves \eqref{eq:actualsurvivors}.

This use of the actual image is decisive. There is no assumption that one may prescribe unrelated Frobenius data independently in a product of groups. The required element is already present: it is the image of one actual complex conjugation. Chebotarev supplies prime Frobenius classes matching that element in the finite quotient.

For $M=1$, the representation product is empty. The cyclotomic quotient modulo $5$ suffices. Finitely many ramified primes are omitted in every case, which does not affect positive density.

\figplaceholder{actual_survivors}{The simultaneous congruence proof uses the actual image of a joint representation. Complex conjugation supplies one existing trace-zero element with cyclotomic value $-1$, and Chebotarev supplies matching prime classes. No independence of the representation factors is assumed.}

\section{All primitive exponents vanish modulo the fixed modulus}
At a prime in \eqref{eq:actualsurvivors}, one has $\chifive(p)=1$, $p^{10}\equiv1\pmod M$, $p^{11}\equiv-1\pmod M$ and $\tau(p)\equiv0\pmod M$. In the formal power series ring $(\Z/M\Z)[[X]]$, the primitive factor becomes
\[
\sum_{r\ge0}R(p^r)X^r
=\frac{1-p^{10}X^2}{1-\tau(p)X+p^{11}X^2}
\equiv\frac{1-X^2}{1-X^2}=1.
\]
The denominator has constant coefficient $1$, so it is invertible even when $M$ is composite or even. No field hypothesis is needed.

Coefficient extraction gives \eqref{eq:allrcongruence}. The constant coefficient remains $R(1)=1$. Equivalently, the reduced Hecke recurrence has $T_0=1$, $T_1=0$ and $T_r=T_{r-2}$. Even $T_r$ are $1$, odd $T_r$ are $0$, and the primitive subtraction kills every positive exponent.

This includes nonzero even primitive blocks. Theorem~\ref{thm:primitiveeven} says they are nonzero integers. The present theorem says that at suitable varying primes they are divisible by any preassigned $M$. The two statements are compatible.

\section{The quantifiers cannot be reversed}
Given a fixed finite list of moduli $m_1,\ldots,m_s$, apply Theorem~\ref{thm:actualsurvivors} to their least common multiple. Infinitely many actual primes defeat every test in that list whose decisive step is finding a nonzero residue of $\tau(p)$ or any $R(p^r)$. This remains true even if the test is allowed all positive local exponents.

The assertion is
\[
\forall M\ \exists\text{ a positive-density set of suitable }p.
\]
It is not
\[
\exists p\ \forall M:\ M\mid\tau(p).
\]
The surviving primes may change with $M$. No passage to an exact integer zero is available by reversing the quantifiers.

Nor does the theorem exclude a modulus chosen adaptively from $p$, such as the $p^6$ criterion. It excludes a fixed finite collection of nonzero-residue tests of the stated type. Congruences used for other purposes and global arguments remain possible.

The fresh arithmetic replay searched twenty-four selected moduli through $10000$. Twenty-two had a simultaneous witness, producing $704$ checked primitive congruences through exponent $32$. The moduli $11$ and $19$ had no witness in that bounded range. Their absence is consistent with a positive-density theorem and does not refute it. Synthetic recurrence tests add $24576$ exact checks, but the universal prime-existence proof is the Galois argument above, not those counts.

\figplaceholder{quantifiers}{For each fixed modulus there are actual prime survivors. The primes may vary with the modulus. The quantifiers therefore do not assert one prime whose integer coefficient is divisible by every modulus.}

\chapter{A locally compatible trace-zero object}\label{ch:local}
\section{The local conditions do not contradict zero trace}
A zero coefficient at an actual prime would mean zero Frobenius trace. It is tempting to expect the weight, slopes or filtration to make that impossible. The following explicit local construction shows that the standard conditions admit such a trace.

\begin{theorem}\label{thm:localphi}
For every prime $p$, there is a weakly admissible two-dimensional filtered $\varphi$-module over $\Q_p$, with monodromy $N=0$, filtration jumps $0,11$ and Frobenius characteristic polynomial $X^2+p^{11}$.
\end{theorem}

This is a local existence theorem. It neither identifies the object with the restriction of $\Delta$'s global representation nor asserts the existence of a level-one modular form with a prescribed zero coefficient.

\section{The matrix and filtration}
Take $D=\Q_pe_1+\Q_pe_2$ and define
\[
\varphi(e_1)=e_2,\qquad\varphi(e_2)=-p^{11}e_1,\qquad N=0.
\]
In the displayed basis,
\[
\varphi=\begin{pmatrix}0&-p^{11}\\1&0\end{pmatrix},\qquad
\Tr\varphi=0,\qquad\det\varphi=p^{11},\qquad
\varphi^2=-p^{11}I.
\]
Choose any $\Q_p$-line $L\subset D$ and put
\[
\Fil^iD=
\begin{cases}
D&i\le0,\\
L&1\le i\le11,\\
0&i\ge12.
\end{cases}
\]
The two nonzero associated graded pieces are $\operatorname{gr}^0D=D/L$ and $\operatorname{gr}^{11}D=L$, each of dimension $1$.

\section{Checking weak admissibility}
For the chosen cohomological convention, define
\[
t_H(D)=\sum_i i\dim\operatorname{gr}^iD,\qquad
t_N(D)=v_p(\det\varphi).
\]
Weak admissibility requires $t_H(D)=t_N(D)$ and $t_H(D')\le t_N(D')$ for every $\varphi$-stable subobject $D'$, with the induced filtration. The total equality is
\[
t_H(D)=11=t_N(D).
\]

There is no proper nonzero $\varphi$-stable $\Q_p$-subspace. In dimension two such a subspace would be a line, and $\varphi$ on it would have an eigenvalue $\lambda\in\Q_p$ satisfying
\[
\lambda^2=-p^{11}.
\]
This would require $2v_p(\lambda)=11$, impossible for an integer-valued valuation on $\Q_p^\times$. Thus the proper-subobject inequalities are vacuous. The zero object imposes equality $0=0$, and the whole object satisfies the total equality. The filtered module is weakly admissible.

The proof works at every $p$, including $2$. It does not rely on a sign argument over $\Q_p$ or on quadratic-residue classifications. The odd determinant valuation is sufficient.

\section{The realization theorem and conventions}
Colmez--Fontaine's theorem makes a weakly admissible filtered module admissible \cite{CF}. In the present setting $K=K_0=\Q_p$ and the base Frobenius is the identity. The theorem supplies a two-dimensional semistable $\Q_p$-representation realizing the filtered module. Since $N=0$, the comparison is crystalline, using $B_{\mathrm{cris}}=B_{\mathrm{st}}^{N=0}$.

The filtration jumps $0,11$ and determinant $p^{11}$ are stated in the cohomological convention. Dualizing changes the Galois and Hodge sign conventions. The arithmetic-Frobenius representation in Chapter~\ref{ch:congruences} and this local filtered object should not be identified merely because a polynomial has been written in a familiar form. The construction's role is compatibility, not that missing identification.

\section{Purity and slopes}
The Newton slopes are $11/2,11/2$. The complex roots of $X^2+p^{11}$ are
\[
\pm i\,p^{11/2}.
\]
They are algebraic integers, since they solve a monic integral polynomial, and each has the expected weight-$11$ absolute value. Thus zero trace is compatible simultaneously with integrality, purity, Deligne's size bound and weak admissibility.

The Newton polygon lies above the Hodge polygon with jumps $0,11$, and the total endpoints agree. There is no contradiction in those polygons. A local condition that is satisfied by this object cannot, by itself, forbid zero trace.

For the actual nonzero $\tau(p)$ at $p\ge5$, Deligne's estimate forces $v_p(\tau(p))\le5$. A value with valuation at least $6$ would be zero by the size bound. This does not exclude the zero case. It is precisely the case the conjecture asks us to eliminate.

\figplaceholder{trace_phi}{The constructed rank-two filtered Frobenius module has trace zero, determinant $p^{11}$ and filtration jumps $0,11$. An odd valuation excludes a rational eigenline, making the proper-subobject inequalities vacuous. This local compatibility is distinct from an actual global identification with $\Delta$.}

\section{What is missing globally}
The actual coefficients belong to one global modular form, with compatible Galois representations and the packet realization already described. The local module constructed here is not asserted to belong to that system. A successful local-to-global argument must use a further property that distinguishes the actual restriction from this compatible object.

Finite matrix checks can authenticate the formula $\varphi^2=-p^{11}I$. The release replay checks it at all $1229$ primes through $10000$. The absence of a rational eigenline at every prime follows from the valuation proof, not from that finite matrix range. Conversely, neither the finite checks nor the realization theorem provides the additional global exclusion needed for Lehmer.

\chapter{A global arithmetic countermodel to the retained premises}\label{ch:model}
\section{An exact test of what the premises imply}
The recurrence, parity, congruences, size bounds and primitive square theorem are all strong facts. It is natural to ask whether their combination already forces nonvanishing. The next construction gives an arithmetic sequence satisfying all of them and having a zero.

\begin{theorem}\label{thm:countermodel}
For every real $B\ge1$ and integer $M\ge1$, there exist a prime $q>B$ and a multiplicative integer sequence $f$ with the following properties:
\begin{enumerate}
\item $f(1)=1$ and $f(n)=\tau(n)$ for every $n\le B$.
\item $f(n)\equiv\tau(n)\pmod{10M}$ for every positive integer $n$.
\item Every weight-$12$ prime-power recurrence holds:
\[
f(p^r)=f(p)f(p^{r-1})-p^{11}f(p^{r-2}),\qquad r\ge2.
\]
\item $|f(n)|\le d(n)n^{11/2}$, where $d(n)$ counts the positive divisors of $n$.
\item $f(n)$ is odd exactly when $n$ is an odd square, and $f(n)\equiv n\sigma(n)\pmod5$.
\item The arithmetic primitive transform $R_f$ is nonzero at every square.
\item Nevertheless $f(q)=R_f(q)=0$.
\end{enumerate}
\end{theorem}

This is a countermodel to an implication from the listed premises alone. It is not a counterexample to Lehmer for the fixed actual form. The sequence is not unconditionally described as nonmodular, because if the chosen actual $\tau(q)$ were already zero, the construction could coincide with $\tau$.

\section{Choosing an actual prime}
Apply Theorem~\ref{thm:actualsurvivors} with $L=10M$. Its positive-density set contains a prime $q>B$ with
\[
q\equiv-1\pmod{50M},\qquad\tau(q)\equiv0\pmod{10M}.
\]
Enlarging $q$ avoids the finitely many prime divisors of $10M$. In particular $q$ is odd, $q\ne5$ and $\chifive(q)=1$.

Set prime parameters
\[
a_p=\tau(p)\quad(p\ne q),\qquad a_q=0.
\]
For every prime define integral local blocks
\[
F_{p,0}=1,\qquad F_{p,1}=a_p,\qquad
F_{p,r}=a_pF_{p,r-1}-p^{11}F_{p,r-2}.
\]
For $n=\prod p^{r_p}$, define $f(n)=\prod_pF_{p,r_p}$. Unique factorization makes this a well-defined multiplicative integer sequence with the stated recurrence.

\section{Agreement at every fixed modulus and finite prefix}
Only the parameter at $q$ has changed. Every $n\le B$ is coprime to $q$, so $f(n)=\tau(n)$ there. At $q$, the old and new first parameters agree modulo $10M$. Induction in the integral recurrence gives
\[
F_{q,r}\equiv\tau(q^r)\pmod{10M}
\]
for every exponent. All other local blocks agree exactly. Multiplying local congruences gives $f(n)\equiv\tau(n)\pmod{10M}$ for every $n$, not only through a finite range.

The parity criterion follows from agreement modulo $2$, and the modulo-$5$ congruence follows from agreement modulo $5$. Thus even globally quantified versions of these retained properties do not remove the zero parameter.

One can incorporate any finite list of fixed congruence moduli by taking their common multiple in $M$. One can incorporate any finite coefficient check by increasing $B$. The resulting sequence and prime may change with $B$ and $M$. No statement that one altered sequence agrees with $\tau$ modulo every modulus simultaneously is made.

\section{The full size bound survives}
At primes other than $q$, the Hecke roots have absolute value $p^{11/2}$ by Deligne. Formula \eqref{eq:rootrec} therefore gives
\[
|\tau(p^r)|\le(r+1)p^{11r/2}.
\]
At the modified prime,
\[
F_{q,2m}=(-q^{11})^m,\qquad F_{q,2m+1}=0.
\]
These obey the same local bound at every exponent, including zero. Multiplying the local estimates gives
\[
|f(n)|\le\prod_{p^{r_p}\Vert n}(r_p+1)p^{11r_p/2}
=d(n)n^{11/2}.
\]
A zero coefficient is fully compatible with an upper bound. The global estimate has not detected the modified parameter.

\section{All primitive squares still survive}
Define the arithmetic transform
\[
R_f(n)=\sum_{d^2\mid n}\mu(d)\chifive(d)d^{10}f(n/d^2).
\]
It is multiplicative. Its prime blocks are $R_f(p)=a_p$, and for $r\ge2$,
\[
R_f(p^r)=F_{p,r}-\chifive(p)p^{10}F_{p,r-2}.
\]
For $p\ne q$, these are the actual $R$ blocks and their even terms are nonzero. At $q$, direct substitution and $\chifive(q)=1$ give
\[
R_f(q^{2m})=(-1)^mq^{11m-1}(q+1)\ne0,\qquad m\ge1.
\]
Thus every positive-exponent factor of $R_f(s^2)$ is nonzero, and $R_f(1)=1$, so the transform is nonzero at every square. Meanwhile $R_f(q)=a_q=0$.

The square property is universal in this construction. It is not merely checked through a chosen exponent. This makes the countermodel stronger than a finite imitation of the actual sequence.

\section{The logical conclusion and the remaining information}
The theorem refutes
\[
\text{the listed arithmetic properties}\quad\Longrightarrow\quad
\text{every coefficient is nonzero}.
\]
The construction does not impose a modular-form realization, the packet-shell identity or the compatible global Galois representations of $\Delta$. Those are potential sources of additional information, rather than properties already contained in the weaker premises.

There is no need to know whether $\tau(q)$ is nonzero. If it is nonzero, $f$ differs from $\tau$. If it is zero, $f=\tau$ and Lehmer is already false. In either case, an arithmetic sequence with the listed properties and a zero exists. Calling $f$ nonmodular without the extra nonzero assumption would be logically stronger than what has been proved.

Any successful resumption must therefore identify a particular additional consequence of the actual global modular form or its actual packet geometry, and show how it excludes a zero prime parameter. More fixed congruences, a larger finite prefix, the same recurrence, the same upper bound and the same all-square theorem do not supply that implication, even together.
